\section{An initial triple of pressure slopes}
\label{i:section}

The narrow zero-energy shell in the core permits clusters of four particles,
with relative volume of order \(l_t^9=\ee^{-9t}\).  This entropy cost,
together with the packing density \(p\), produces three competing branches
after the first ideal attraction.  We show that, for one sufficiently large
finite \(t\), the pressure subdifferential throughout a parameter neighborhood
is generated by actual supports near the three limiting slopes.  The
subsequent iteration will preserve this separation and find a parameter
where all three groups are present.

In this section alone we let \(t\) tend to infinity to obtain the required
estimates; we fix its finite value at the end.

\subsection{Clusters and the bare-core pressure}

\begin{proposition}[Pressure of the bare core]
\label{i:bare-pressure}
Let \(q^{(t)}=Q_{\psi_t}\), with the pressure normalization and the core
of \Cref{f:thermodynamics,f:core}.  Then, locally uniformly on
\(\Omega=(0,\infty)\times\R\),
\begin{equation}
 q^{(t)}(\beta,h)\longrightarrow
 q_*(\beta,h):=\max\{0,ph,4ph-9p\}.
 \label{i:bare-limit}
\end{equation}
\end{proposition}

\begin{proof}
Fix a compact parameter set \(K\subset\Omega\), and choose
\(\beta_->0\) and \(H<\infty\) such that
\(\beta\ge\beta_-\) and \(\abs{h}\le H\) on \(K\).
Throughout the proof \(V=L^3\), and
\(I_t=[l_t,(1+\Delta)l_t]\).

\emph{Free configurations and the upper bound.}
Call a pair bad when its distance is less than one and is not in \(I_t\).
A configuration is free if it has no bad pair.  In a free configuration,
join distinct particles whose distance is less than one.  Every edge then
has length at most \((1+\Delta)l_t\).  The endpoints of a path of two
edges are therefore at distance at most \(2(1+\Delta)l_t<1\), and hence
are joined.  Each connected component is consequently a clique with all
pair distances in \(I_t\).  By \Cref{f:shell-graph}, its size is at most
four.  Choosing one representative
from every component gives a set with pairwise distances at least one,
so the component count is at most \(M_L\).

Consider the grand integral over free configurations with their energy
discarded:
\[
 \Xi_{L,\mathrm{free}}(h)
 =\sum_{n\ge0}\frac{\ee^{thn}}{n!}
       \int_{\Lambda_L^n}\mathbf 1_{\mathrm{free}}(x_1,\ldots,x_n)
          \dd x_1\cdots\dd x_n.
\]
For an ordered list of component sizes \(m_1,\ldots,m_k\), with
\(m_i\in\{1,2,3,4\}\) and \(\sum_i m_i=n\), the number of ordered
partitions of the particle labels into these blocks is
\(n!/\prod_i m_i!\).  On summing over all ordered size lists, division
by \(k!\) counts each unordered family of label blocks once.  This
remains true when sizes repeat: the label blocks themselves are distinct
and have exactly \(k!\) orderings.

In a block of size \(m\), choose its least label as the root.  Its position has volume
at most \(V\), and each of the other \(m-1\) particles lies in a ball of
volume \(C l_t^3\) about it, where \(C=4\pi(1+\Delta)^3/3\).
After cancellation of the global factor \(1/n!\), the block therefore
contributes at most \(V(C l_t^3)^{m-1}/m!\).  Discarding all further
restrictions within and between components and summing their sizes gives
\begin{equation}
 \Xi_{L,\mathrm{free}}(h)
 \le\sum_{k=0}^{M_L}\frac{(V\Theta_t(h))^k}{k!},
 \qquad
 \Theta_t(h)=\sum_{m=1}^4
       \frac{\ee^{thm}(C l_t^3)^{m-1}}{m!}.
 \label{i:rooted-count}
\end{equation}

To pass from free configurations to all configurations, let \(B(X)\)
be the set of labels incident to a bad pair.  There are at least
\(\abs{B(X)}/2\) bad pairs, and each has core energy exactly \(t^2\).
All other core terms are nonnegative, and \(X\setminus B(X)\) is free.
Thus the following pointwise bound holds:
\begin{equation}
 \ee^{-\beta U^{\psi_t}(X)}
 \le\sum_{J\subset\{1,\ldots,n\}}
       \ee^{-\beta t^2\abs{J}/2}
       \mathbf 1_{\mathrm{free}}(X\setminus J).
 \label{i:removal-pointwise}
\end{equation}
Indeed, the summand \(J=B(X)\) already bounds the left side.
For \(r=\abs{J}\) and \(m=n-r\), the label factor is
\(\binom{m+r}{r}/(m+r)!=1/(r!m!)\).  The removed positions integrate
to \(V^r\).  Tonelli's theorem and the grand sum now give the exact
factorization of the upper bound
\begin{equation}
 \Xi_L^{\psi_t}(\beta,h)
 \le
 \exp\!\left(V\ee^{th-\beta t^2/2}\right)
       \Xi_{L,\mathrm{free}}(h).
 \label{i:removal-factor}
\end{equation}
Shell endpoints are included in \(I_t\), although their energy may be
positive.  Pairs at distance exactly one are not edges of the component
graph.  Both conventions merely discard nonnegative energy in the upper
bound, so these arguments hold pointwise, including at the endpoints.

Set \(a_t(h)=t^{-1}\log\Theta_t(h)\).  The four exponential rates in
\(\Theta_t\) are
\[
 h,\qquad 2h-3,\qquad 3h-6,\qquad 4h-9.
\]
The middle two are convex combinations of the first and last.  Their
positive prefactors in \eqref{i:rooted-count} are fixed, so
\begin{equation}
 a_t(h)=\max\{h,4h-9\}+O(t^{-1}),
 \label{i:cluster-rates}
\end{equation}
uniformly for \(\abs{h}\le H\).  For \(k\ge1\), the inequality
\(k!\ge(k/\ee)^k\), with \(x=k/V\), gives
\[
 \frac1{tV}\log\frac{(V\Theta_t)^k}{k!}
 \le x a_t(h)+\frac{x(1-\log x)}t
 \le \frac{M_L}{V}\max\{0,a_t(h)\}+\frac1t,
\]
since \(x(1-\log x)\le1\) for \(x>0\).  The same final bound covers
\(k=0\).  There are only \(M_L+1=O(1+V)\) terms in
\eqref{i:rooted-count}.  First letting \(L\to\infty\), with \(t\)
fixed, and using \Cref{f:packing,f:thermodynamics}, we obtain
\begin{equation}
 q^{(t)}(\beta,h)
 \le p\max\{0,a_t(h)\}+\frac1t
       +\frac1t\ee^{th-\beta t^2/2}.
 \label{i:upper-pressure}
\end{equation}
The final term is at most
\(t^{-1}\ee^{tH-\beta_-t^2/2}\) on \(K\).  Thus
\eqref{i:cluster-rates} proves the required uniform upper bound.

\emph{Zero-energy configurations and the lower bound.}
Fix \(s>0\) and \(0<\nu\le1\).  Dilate a maximizing packing in a
side-\(s\) cube by \(1+\nu\), and center it in a cube of side
\[
 P=(s+2)(1+\nu).
\]
Repeat this pattern in adjacent period cubes.  Centers in one period
have separation at least \(1+\nu\); centers in different periods have
at least this separation as well.  Every center has distance at least
\(1+\nu\) from the boundary of its period cube.  A cube of side \(mP\)
contains exactly \(k=m^3M_s\) centers.

Choose \(r\in(0,\nu/16)\) and let \(b_0=4\pi r^3/3\).
At every center \(z\), put the root particle in \(B(z,r)\).
There are two constructions, with cluster size \(j=1\) or \(j=4\).
For \(j=4\), choose vectors \(a_2,a_3,a_4\in\R^3\) such that
\(0,a_2,a_3,a_4\) are the vertices of a regular tetrahedron of side
\(1+\Delta/2\).  Fix \(\delta\in(0,\Delta/24)\) and require
\[
 x_i-x_1\in B(l_t a_i,\delta l_t),\qquad i=2,3,4.
\]
A root-to-vertex distance differs from
\((1+\Delta/2)l_t\) by at most \(\delta l_t\), and another edge
differs by at most \(2\delta l_t\).  All six distances consequently
lie strictly inside
\[
 ((1+\Delta/3)l_t,(1+2\Delta/3)l_t),
\]
where the core vanishes.  Write \(b_1=4\pi\delta^3/3>0\).
The change from positions to the root and the three relative vectors
has Jacobian one.  For one prescribed label ordering the integration
volume at each center is therefore
\(b_0(b_1l_t^3)^3\), with positive relative volume \(b_1^3l_t^9\).
For \(j=1\), the corresponding volume is \(b_0\).

Every particle lies within \(r+2l_t\) of its center.  For all sufficiently
large \(t\), depending on \(\nu,r\) but not on \(m,\beta,h\), these
center neighborhoods are disjoint and lie inside their period cubes,
and intercluster distances are greater than
\(1+\nu/2>1+1/t\).  Hence all configurations just described have
exactly zero core energy.

The centers are distinguishable.  Assign an unordered set of \(j\)
particle labels to each of the \(k\) centers, and retain just one
deterministic ordering of each assigned set, for example the increasing
order.  There are \((jk)!/(j!)^k\) such assignments.  Different
assignments send some label to different disjoint center neighborhoods,
so their integration events are disjoint.  Possible alternative choices
of the root within the same set are never summed.  Dividing by the
canonical factorial \((jk)!\) gives
\begin{equation}
 Z^{\psi_t}_{mP,jk}(\beta)
 \ge\left(\frac{b_0(b_1l_t^3)^{j-1}}{j!}\right)^k,
 \qquad j\in\{1,4\}.
 \label{i:cluster-lower}
\end{equation}
Using this term in the grand sum and first letting \(m\to\infty\), we
find
\begin{equation}
 q^{(t)}(\beta,h)
 \ge c_{s,\nu}\left(jh-3(j-1)
       +\frac1t\log\frac{b_0b_1^{j-1}}{j!}\right),
 \qquad
 c_{s,\nu}=\frac{M_s}{((s+2)(1+\nu))^3}.
 \label{i:lower-pressure}
\end{equation}
The vacuum also gives \(q^{(t)}\ge0\).  The constants in
\eqref{i:lower-pressure} have no dependence on \(\beta,h\), and
\(c_{s,\nu}\) can be made arbitrarily close to \(p\) by taking \(s\)
large and \(\nu\) small.  More explicitly, on \(\abs{h}\le H\) the
two functions \(jh-3(j-1)\), \(j=1,4\), have a common finite absolute
bound.  Given any error tolerance, first choose \(s,\nu\) so that
replacing \(c_{s,\nu}\) by \(p\) changes both branches by less than
that tolerance, and then take \(t\) large enough to make the two
logarithmic terms small.  Taking the maximum of these branches and zero
preserves the uniform error, even where a nonvacuum branch is negative.
Together with \eqref{i:upper-pressure}, this proves
\eqref{i:bare-limit} on \(K\).
\end{proof}

\subsection{The limiting attraction and its active densities}

We next apply the ideal attraction with mass \(A=3/(2p)\).  Let
\(P_1^t=T_Aq^{(t)}\), and define
\begin{equation}
 \begin{split}
 L_1(\beta,h)&=0,\\
 L_2(\beta,h)&=ph+\frac{\beta A p^2}{2},\\
 L_3(\beta,h)&=4ph-9p+\frac{\beta A(4p)^2}{2},\\
 P_*(\beta,h)&=\max_{1\le i\le3}L_i(\beta,h),
 \qquad D_*=\{0,p,4p\}.
 \end{split}
 \label{i:three-planes}
\end{equation}

\begin{lemma}[The limiting transform and its optimizer gap]
\label{i:limiting-transform}
The functions \(P_1^t\) converge locally uniformly on \(\Omega\) to
\(P_*\).  For every compact \(K\subset\Omega\), all their active
optimizers lie in one interval \([0,U]\) for all sufficiently large
\(t\), and
\[
 \sup_{\lambda\in K}\ \sup_{u\ \mathrm{active\ for}\ P_1^t(\lambda)}
       \dist(u,D_*)\longrightarrow0.
\]
In particular, this holds at parameters where several limiting planes
meet.  If \(\beta_-=\min_{(\beta,h)\in K}\beta\), the limiting
objective obeys the explicit bound
\begin{equation}
 F_*(\lambda,u):=q_*(\beta,h+\beta Au)-\frac{\beta A u^2}{2}
 \le P_*(\lambda)-\frac{\beta_-A}{2}\dist(u,D_*)^2
 \quad(\lambda\in K,\ u\in\R).
 \label{i:objective-gap}
\end{equation}
\end{lemma}

\begin{proof}
Apply the uniform-in-\(t\) bound in \Cref{k:transform-properties} to
the bare core \(\psi_t\), mass \(A\), and parameter compact \(K\).
For all sufficiently large \(t\), every active shift lies in one fixed
interval \([0,U]\).  Enlarge \(U\), if necessary, so that \(4p\le U\).

Write \(d_1=0,d_2=p,d_3=4p\).  Completing the square gives the identity
\begin{equation}
 F_*(\lambda,u)
 =\max_{1\le i\le3}
       \left\{L_i(\lambda)-\frac{\beta A}{2}(u-d_i)^2\right\}.
 \label{i:limiting-parabolas}
\end{equation}
Its supremum over \(u\) is \(P_*(\lambda)\), and all its maximizing
values belong to \(D_*\).  The same identity proves
\eqref{i:objective-gap}.

The shifted points
\[
 K'=\{(\beta,h+\beta Au):(\beta,h)\in K,\ 0\le u\le U\}
\]
form a compact subset of \(\Omega\).  Put
\(E_t=\sup_{K'}\abs{q^{(t)}-q_*}\), which tends to zero by
\Cref{i:bare-pressure}.  Both transforms can be maximized over
\([0,U]\), so
\[
 \sup_K\abs{P_1^t-P_*}\le E_t.
\]
If \(u\) is active for \(P_1^t(\lambda)\), approximation of its
objective and of the maximum shows
\(F_*(\lambda,u)\ge P_*(\lambda)-2E_t\).  By
\eqref{i:objective-gap},
\begin{equation}
 \dist(u,D_*)
 \le 2\sqrt{\frac{E_t}{\beta_-A}}.
 \label{i:active-density-error}
\end{equation}
This estimate is uniform on \(K\) and requires no strict inequality
between the three affine branches.
\end{proof}

\subsection{Uniform separation of generating slopes}

The optimizer gap locates the density coordinates of active supports.
We now identify neighborhoods of the full slopes that will remain separated
under the iteration.  The slopes of the three limiting planes and their
barycenter are
\begin{equation}
 \begin{gathered}
 d_1=0,\qquad d_2=p,\qquad d_3=4p,\qquad
 v_i=(Ad_i^2/2,d_i),\\
 v_1=(0,0),\qquad v_2=(3p/4,p),\qquad v_3=(12p,4p),\\
 g=\frac{v_1+v_2+v_3}{3}=(17p/4,5p/3),\qquad
 \lambda_*=(1,-3/4).
 \end{gathered}
 \label{i:slopes}
\end{equation}
All three planes vanish at \(\lambda_*\), and
\(\det(v_2-v_1,v_3-v_1)=-9p^2\ne0\).  Thus \(g\) is strictly
inside their slope triangle.

\begin{lemma}[Robust slope geometry]
\label{i:geometry}
There are constants \(0<\tau<p/16\) and \(\gamma>0\) with the
following properties.  Every choice
\(w_i\in\clB{v_i}{\tau}\), \(i=1,2,3\), satisfies
\begin{equation}
 \clB{g}{\gamma}\subset\conv\{w_1,w_2,w_3\}.
 \label{i:robust-triangle}
\end{equation}
For each distinct pair \(i,j\),
\begin{equation}
 g\notin\conv\bigl(\clB{v_i}{\tau}\cup\clB{v_j}{\tau}\bigr).
 \label{i:two-groups}
\end{equation}
Moreover, for every \(\lambda\in\Omega\),
\begin{equation}
 P_*(\lambda)-P_*(\lambda_*)
       -\ip{g}{\lambda-\lambda_*}
 \ge\gamma\norm{\lambda-\lambda_*}.
 \label{i:linear-growth}
\end{equation}
\end{lemma}

\begin{proof}
Because \(g\) is an interior point of the nondegenerate triangle, the
two numbers
\[
 \eta=\min_{\norm{z}=1}\max_i\ip{v_i-g}{z}>0,
 \qquad
 \zeta=\min_{i<j}\dist\bigl(g,[v_i,v_j]\bigr)>0
\]
are positive.  For the first assertion, one may also see positivity
directly: if all three inner products were nonpositive, their sum would
be zero, so all would vanish, contradicting noncollinearity.
Choose
\[
 0<\tau<\min\{p/16,\eta/4,\zeta/2\},
 \qquad \gamma=\eta/2.
\]
For any unit vector \(z\) and any allowed vertices \(w_i\),
\[
 \max_i\ip{w_i-g}{z}\ge\eta-\tau>\gamma.
\]
The support-function characterization of inclusion of compact convex
sets proves \eqref{i:robust-triangle}.  Every point of the hull of two
radius-\(\tau\) balls lies within distance \(\tau\) of the segment
joining their centers.  Since \(\tau<\zeta\), this also proves
\eqref{i:two-groups}.  Finally, all three planes agree at
\(\lambda_*\), and hence for \(z=\lambda-\lambda_*\),
\[
 P_*(\lambda)-P_*(\lambda_*)-\ip{g}{z}
 =\max_i\ip{v_i-g}{z}\ge\eta\norm{z}\ge\gamma\norm{z}.
\]
\end{proof}

The geometry of \Cref{i:geometry} is illustrated in
\Cref{fig:support-geometry}.
\begin{figure}[tb]
\centering
\begin{tikzpicture}[>=Latex,font=\small]
  \begin{scope}[scale=.48]
    \draw[->,black!55] (-.45,0) -- (13.1,0)
      node[below=2pt,text=black] {\(e/p\)};
    \draw[->,black!55] (0,-.35) -- (0,5.0)
      node[above,text=black] {\(d/p\)};
    \fill[blue!8] (0,0) -- (.75,1) -- (12,4) -- cycle;
    \draw[blue!60!black,thick] (0,0) -- (.75,1) -- (12,4) -- cycle;
    \fill (0,0) circle (2.4pt) node[below left] {\(v_1\)};
    \fill (.75,1) circle (2.4pt) node[above left] {\(v_2\)};
    \fill (12,4) circle (2.4pt) node[above right] {\(v_3\)};
    \fill[orange!75!black] (4.25,1.666667) circle (2.4pt);
    \draw[orange!85!black,thick] (4.25,1.666667) circle (.115);
    \node[above left] at (4.0,2.25) {\(g\)};
    \draw[black!55] (3.93,2.12) -- (4.18,1.80);
    \draw[black!55,dashed] (3.75,1.15) rectangle (4.75,2.2);
  \end{scope}
  \draw[black!55,->] (2.28,1.05) to[bend left=14] (7.48,2.3);
  \begin{scope}[shift={(8.9,1.65)}]
    \fill[orange!8] (0,0) circle (1.35);
    \draw[orange!85!black,thick] (0,0) circle (1.35);
    \draw[black!55,dashed] (-1.7,0) -- (1.7,0);
    \draw[very thick,blue!65!black] (-.93,0) -- (.93,0);
    \fill[blue!65!black] (-.93,0) circle (2pt)
      node[below=5pt] {\((s_1,\rho)\)};
    \fill[blue!65!black] (.93,0) circle (2pt)
      node[below=5pt] {\((s_2,\rho)\)};
    \fill (0,0) circle (1.6pt) node[above=4pt] {\(g\)};
    \node at (0,1.67) {\(\overline B(g,\gamma)\)};
    \node at (0,-1.72) {\(\rho=5p/3,\quad s_1<s_2\)};
  \end{scope}
\end{tikzpicture}
\caption{Pressure-support coordinates \((e,d)\), with the left panel
drawn in coordinates \((e/p,d/p)\).
The triangle of the initial limiting slopes \(v_i\) contains a ball
about its barycenter \(g\);
the ball persists under the controlled perturbations of its vertices.
The schematic inset shows two points on the central horizontal slice
at density \(5p/3\).  Once the ball lies in the subdifferential of a
pressure at one parameter, each nearby horizontal slice supplies two
supports with distinct inverse-temperature slopes at that same parameter.
Lemma~\ref{m:canonical-corner} converts this geometry into the canonical
derivative jump throughout a density interval.
}
\label{fig:support-geometry}
\end{figure}


Define the closed parameter neighborhoods
\begin{equation}
 W(r)=\clB{\lambda_*}{r},\qquad W_1=W(1/4)\subset\Omega.
 \label{i:neighborhoods}
\end{equation}
The next proposition gives a statement about actual supporting slopes
throughout \(W_1\), as required for the subsequent iteration.

\begin{proposition}[Initial separation of generating slopes]
\label{i:initial-separation}
For all sufficiently large finite \(t\), the size conditions in
\Cref{f:stability,k:transform-properties} hold and
\begin{equation}
 \sup_{\lambda\in W_1}\abs{P_1^t(\lambda)-P_*(\lambda)}
       <\frac{\gamma}{64}.
 \label{i:initial-pressure-error}
\end{equation}
At every \(\lambda\in W_1\) one also has the exact identity
\begin{equation}
 \partial P_1^t(\lambda)
 =\conv\left(
       \partial P_1^t(\lambda)\cap
       \bigcup_{i=1}^3\clB{v_i}{\tau/2}\right).
 \label{i:initial-hull}
\end{equation}
Fix one such \(t\) for the rest of the construction.
\end{proposition}

\begin{proof}
Use the uniform optimizer bound \([0,U]\) from
\Cref{i:limiting-transform} with \(K=W_1\).  All shifted points
\(x=(\beta,h+\beta Au)\) lie in the fixed compact \(K'\) from that
proof.  Choose one \(\sigma>0\) such that
\[
 K''=\{x+(s,0):x\in K',\ \abs{s}\le\sigma\}\subset\Omega,
\]
and set \(E'_t=\sup_{K''}\abs{q^{(t)}-q_*}\to0\).
For every \((e,d)\in\partial q^{(t)}(x)\), its two supporting
inequalities at \(x+(\sigma,0)\) and \(x-(\sigma,0)\), together
with the independence of \(q_*\) from \(\beta\), imply
\begin{equation}
 -\frac{2E'_t}{\sigma}\le e\le\frac{2E'_t}{\sigma}.
 \label{i:energy-support-error}
\end{equation}
These are fixed steps on an enlarged compact, so the estimate also
covers outer parameters on the boundary of \(W_1\).

At a shifted point corresponding to an active \(u\),
\Cref{k:active-parabola} gives \(d=u\) for every underlying support.
By \Cref{k:subdifferential}, the actual active generators of
\(\partial P_1^t(\lambda)\) are precisely the vectors in the union
\[
 G_t(\lambda)=
 \bigcup_{u\ \mathrm{active}}
   \left\{(e+Au^2/2,u):(e,u)\in
           \partial q^{(t)}(\beta,h+\beta Au)\right\},
 \qquad
 \partial P_1^t(\lambda)=\conv G_t(\lambda).
\]
For a nearest \(d_i\in D_*\), every such generator \(w\)
satisfies
\[
 \norm{w-v_i}
 \le\frac{2E'_t}{\sigma}
       +\left(1+\frac A2(U+4p)\right)\abs{u-d_i}.
\]
Equations \eqref{i:active-density-error} and
\eqref{i:energy-support-error} show that this tends uniformly to zero
over all outer parameters, all active optimizers, and all underlying
supports.  Thus, for all sufficiently large \(t\),
\[
 G_t(\lambda)\subset
       \partial P_1^t(\lambda)\cap
       \bigcup_i\clB{v_i}{\tau/2}
       \qquad(\lambda\in W_1).
\]
Taking convex hulls proves one inclusion in \eqref{i:initial-hull};
the other follows from convexity of the subdifferential.  Finally,
\eqref{i:initial-pressure-error} follows from
\Cref{i:limiting-transform}.  All earlier size conditions are lower
bounds on \(t\), so they can be imposed at the same time.
\end{proof}
